
\documentclass{article}
\usepackage{colortbl}
\usepackage{makecell}
\usepackage{multirow}
\usepackage{supertabular}

\begin{document}

\newcounter{utterance}

\centering \large Interaction Transcript for game `cladder', experiment `full\_v1.5\_default', episode 5038 with reflexion{-}Qwen3.5{-}4B.
\vspace{24pt}

{ \footnotesize  \setcounter{utterance}{1}
\setlength{\tabcolsep}{0pt}
\begin{supertabular}{c@{$\;$}|p{.15\linewidth}@{}p{.15\linewidth}p{.15\linewidth}p{.15\linewidth}p{.15\linewidth}p{.15\linewidth}}
    \# & \multicolumn{2}{c}{Player} && \multicolumn{2}{c}{Game Master} \\
    \hline

    \theutterance \stepcounter{utterance}  
    & & & \multicolumn{4}{p{0.6\linewidth}}{
        \cellcolor[rgb]{0.9,0.9,0.9}{
            \makecell[{{p{\linewidth}}}]{
                \texttt{\tiny{[P1$\langle$GM]}}
                \texttt{You are an expert in causal inference. The following question is not a typical commonsense query, but rather a meticulously designed question created by a professor specializing in causal inference, intended to assess the students' mastery of the course content.} \\
\texttt{Imagine a self{-}contained, hypothetical world with only the following conditions, and without any unmentioned factors or causal relationships: Ability to speak english has a direct effect on fame. Talent has a direct effect on fame. The overall probability of speaking english is 89\%. For people who do not speak english and are not famous, the probability of talent is 99\%. For people who do not speak english and are famous, the probability of talent is 97\%. For people who speak english and are not famous, the probability of talent is 99\%. For people who speak english and are famous, the probability of talent is 57\%. If we look at people who are famous, does the chance of talent decrease when speaking english?} \\
\texttt{Start your answer with 'Yes' or 'No', followed by additional reasoning or evidence to support your explanation.} \\
            }
        }
    }
     \\ \\

    \theutterance \stepcounter{utterance}  
    & \multicolumn{4}{p{0.6\linewidth}}{
        \cellcolor[rgb]{0.561,0.737,0.561}{
            \makecell[{{p{\linewidth}}}]{
                \texttt{\tiny{[P1$\rangle$GM]}}
                \texttt{Yes} \\
\\ 
\texttt{**Reasoning:**} \\
\\ 
\texttt{To determine whether the chance of talent decreases when speaking English among famous people, we must compare the conditional probability of talent given fame and English (\$P(\textbackslash{}text\{Talent\} \textbackslash{}mid \textbackslash{}text\{Famous\} \textbackslash{}cap \textbackslash{}text\{English\})\$) with the conditional probability of talent given fame and no English (\$P(\textbackslash{}text\{Talent\} \textbackslash{}mid \textbackslash{}text\{Famous\} \textbackslash{}cap \textbackslash{}neg\textbackslash{}text\{English\})\$).} \\
\\ 
\texttt{Based on the conditions provided in the hypothetical world:} \\
\texttt{1.  **For people who do not speak English and are famous:** The probability of talent is **97\%**.} \\
\texttt{2.  **For people who speak English and are famous:** The probability of talent is **57\%**.} \\
\\ 
\texttt{Comparing these two values:} \\
\texttt{\$\$57\textbackslash{}\% < 97\textbackslash{}\%\$\$} \\
\\ 
\texttt{Therefore, within this specific world, the probability of having talent is indeed lower for famous individuals who speak English compared to those who do not. This result illustrates a classic case of **collider bias** (or Berkson's paradox). "Fame" acts as a collider on the path between "Speaking English" and "Talent" (since both cause Fame). Conditioning on the collider (by restricting our observation to only famous people) induces a spurious negative association between the two causes, even though they are independent in the general population.} \\
            }
        }
    }
    & & \\ \\

    \theutterance \stepcounter{utterance}  
    & & & \multicolumn{2}{p{0.3\linewidth}}{
        \cellcolor[rgb]{0.9,0.9,0.9}{
            \makecell[{{p{\linewidth}}}]{
                \texttt{\tiny{[GM$|$GM]}}
                \texttt{Yes} \\
            }
        }
    }
    & & \\ \\

    \theutterance \stepcounter{utterance}  
    & & & \multicolumn{2}{p{0.3\linewidth}}{
        \cellcolor[rgb]{0.9,0.9,0.9}{
            \makecell[{{p{\linewidth}}}]{
                \texttt{\tiny{[GM$|$GM]}}
                \texttt{yes} \\
            }
        }
    }
    & & \\ \\

    \theutterance \stepcounter{utterance}  
    & & & \multicolumn{2}{p{0.3\linewidth}}{
        \cellcolor[rgb]{0.9,0.9,0.9}{
            \makecell[{{p{\linewidth}}}]{
                \texttt{\tiny{[GM$|$GM]}}
                \texttt{game\_result = WIN} \\
            }
        }
    }
    & & \\ \\

\end{supertabular}
}

\end{document}
